\section{Modular labels and information that a norm forgets}\label{sec:labels}
The recovered modular comparison keeps a candidate form, a prime of its coefficient field above seven, and a residue embedding. These labels must remain compatible across auxiliary primes. Taking norms or minimal polynomials can produce a sound necessary condition while forgetting those labels.

Suppose a Hecke eigenvalue $a_{\mathfrak q}$ is replaced by its minimal polynomial. A resultant against a possible curve trace polynomial asks whether some conjugate can match some allowable trace. Repeating the test at a different prime can choose a different conjugate. If the original representation requires one residue embedding for every prime, this relaxed sequence of tests may preserve candidates that the labelled comparison would exclude. An exclusion obtained from the relaxation remains valid. Survival becomes less informative.

The Block 47 audit identified this mechanism in the pinned Pacetti--Villagra Torcomian code. It also identified a parameter dictionary issue at eleven: the unlabelled multisets agreed, but individual labels matched after $t\mapsto t^{-1}$. A refinement attached to the wrong labels could incorrectly combine local constraints. The new reader therefore distinguishes direct point counts on \eqref{eq:sextic} from imported hypergeometric parameter tables.

Two exact local restrictions were established under the stated modular and multiplicative prime comparisons. A rank two trace zero comparison at a prime above eleven forces $11\mid A$. At a prime above nineteen it forces $19\mid B$. In either case the full four dimensional reduction to check is
\[
(T^2+q)^2\pmod7,
\]
which is stronger than zero four dimensional trace. At bad $q$, the multiplicative comparison gives $\pm(q+1)$, nonzero modulo seven for these two primes. At good $q$, the direct curve table selects the stated base zero stratum. These are divisibility restrictions, not contradictions.

\subsection{From point counts to a common polynomial}
For a smooth genus two curve over $\F_q$, let $N_1,N_2$ count its points over $\F_q,\F_{q^2}$, including infinity. Write
\[
s_1=q+1-N_1,\qquad
s_2=\frac{s_1^2-(q^2+1-N_2)}2.
\]
Its Frobenius polynomial is
\begin{equation}\label{eq:quartic}
T^4-s_1T^3+s_2T^2-qs_1T+q^2.
\end{equation}
In the RM setting, the rank two traces are the roots of
\[
X^2-s_1X+(s_2-2q).
\]
The two roots belong to one common quartic. Choosing one root from one fiber and another root from an unrelated fiber can enlarge a necessary locus, but does not give a compatible simultaneous reduction of the original curve. The later Frobenius sieves explicitly retain this common polynomial condition.

\subsection{The remaining modular branches}
Block 48 and the accumulated Block 49 notes retain conditional local type and divisibility reductions, including the branch $3990\mid C$. They do not certify every Hilbert modular candidate as eliminated. The imported authors' transcripts have their own index conventions and some aggregate exception lists do not identify a complete seven specific set. A form index in such a transcript is not an intrinsic arithmetic identifier.

The surviving CM branch also requires care. CM can force a trace to vanish at suitable inert primes, which produces base divisibility constraints such as the one at eleven. It does not imply that every representation with a compatible residual trace is the CM representation itself. A congruence modulo seven is weaker than equality of characteristic zero representations.

The packet interpretation here is exact and limited. A packet consists of a candidate, its coefficient field, the chosen prime above seven, the residue embedding, and the paired local polynomial. Transfer to an unlabelled norm is a projection of that packet. Recovering the forgotten labels is a legitimate strengthening. The required exports were not present in every recovered modular data set, so no completed all form elimination is inferred from the proposed strengthening.
